Two massive stars A and B with masses $m_A$ and $m_B$ are separated by a
distance $d$. Both stars orbit around their center of mass under
gravitational force. Assume their orbits are circular and lie on the
$X$--$Y$ plane whose origin is at the stars' center of mass (see Figure 2).

\begin{description}
\item[(a)] Find the expressions for the tangential and angular speeds of
  star A.
\end{description}

An observer standing on the $Y$--$Z$ plane (see Figure 2) sees the stars from
a large distance with an angle $\theta$ relatively to the $Z$-axis. He
measures that the velocity component of star A to his line of sight has the
form $K\cos(\omega t + \varepsilon)$, where $K$ and $\varepsilon$ are
positive.

\begin{description}
\item[(b)] Express $K^3/\omega G$ in terms of $m_A$, $m_B$, and $\theta$,
  where $G$ is the universal gravitational constant.
\end{description}

Assume that the observer then identifies that star A has mass equal to
$30\,M_S$ where $M_S$ is the Sun's mass. In addition, he observes that star B
produces X-rays and then realizes that it could be a neutron star or a black
hole. This conclusion would depend on $m_B$, i.e.: i) if $m_B < 2M_S$, then B
is a neutron star; ii) if $m_B > 2M_S$, then B is a black hole.

\begin{description}
\item[(c)] A measurement by the observer shows that
  $\dfrac{K^3}{\omega G} = \dfrac{1}{250}\,M_S$. In practice, the value of
  $\theta$ is usually not known. What is the condition on $\theta$ for star B
  to be a black hole?
\end{description}
